\section{Conclusion}
\label{bdc:sec:next}

We have described the BDC process through its cell-fate probabilities and first
two moments. The resulting formulae cover the intracellular load, the released
count, and the first two conditional moments of the burst size, all through the
single coordinate $I(t)$.

The calculation closes because the non-catastrophe probability obeys an
autonomous Riccati equation. The open moment hierarchy is still present, but it
need not be truncated: it can be approached from $I$. The two roots of the
characteristic quadratic then make the closed forms readable, and they retain
their probabilistic meaning at the boundaries recorded in
\cref{bdc:app:limits}.

One reading takes rupture as the lysis of an infected macrophage. The formulae
do not depend on that reading. They do depend on the four assumptions of
\cref{bdc:sec:intro}: linear birth, catastrophe hazard $\delta n$,
instantaneous complete release, and constant rates. Measurement of
$(\lambda,\mu,\delta)$ and the full burst law lie outside this chapter.

\subsection{Completing the single-cell theory}
\label{bdc:sec:next-singlecell}

\fwd{singlecell}{Chapter~5} solves the probability generating function of the
process. This yields the state probabilities $p_n(t)$, rather than only their
sums $I$, $D$ and $\Ifix$, and identifies the quasi-stationary distribution of
the intracellular load conditioned on continued productivity. It also gives the
full burst-time and burst-size laws, including the defective finite-time mass
\[
P_k=\pr{\Kb=k,\ \tau<\infty}
=\frac{\delta}{\lambda}a^{-k},
\qquad k\ge1,
\]
and the conditional laws of size given time and time given size. The several
founder formulae of \cref{bdc:sec:multifounder} then describe cells infected by
more than one particle. Finally, the immediate-transfer model chains successive
cells together by allowing one burst to found the next infection without delay.

What passes from the present chapter to that analysis is deliberately limited.
We know the first two conditional burst moments, and they point towards a
geometric distribution. The generating function supplies the distribution
itself and explains why it agrees with the quasi-stationary load.

\subsection{Embedding the cell in a population}
\label{bdc:sec:next-population}

\fwd{population}{Chapter~6} turns from one infected cell to many. A burst is no
longer an endpoint: released particles may infect fresh cells, each of which
begins its own BDC trajectory. The quantities derived here then become
age-dependent kernels. The survival kernel $\Ifix(\alpha)$ gives the probability
that a cell infected $\alpha$ time units ago remains productive, while
\[
g(\alpha)=\delta K(\alpha)
\]
gives its expected rate of release. Because these functions can be written in
the $I$-coordinate, the resulting renewal equations retain much of the analytic
structure of the single-cell model.

The renewal formulation belongs to the classical theory of age-structured
populations \cite{mckendrick1926applications,vonfoerster1959some}. It also makes
clear what is hidden by the constant-release ordinary differential equations
often used in viral dynamics
\cite{perelson1996hiv,nowak1996population}: release from a
lytic cell has an age profile and occurs in a single burst. In regimes where
that timing matters, a constant-release approximation can be poor. The
comparison with continuous budding release is developed in the population
chapter \cite{pearson2011stochastic,gilchrist2006evolution,wang2006lysis}.

\FloatBarrier
